\documentclass[conference]{IEEEtran}

% ================= PACKAGES =================
\usepackage[utf8]{inputenc}
\usepackage[T1]{fontenc}
\usepackage{amsmath}
\usepackage{amssymb}
\usepackage{graphicx}
\usepackage{xcolor}
\usepackage[hypertexnames=false]{hyperref}
\usepackage{booktabs}
\usepackage{balance}
\hypersetup{
    colorlinks=true,
    urlcolor=blue,
    citecolor=blue,
    linkcolor=black
}

% ================= TITLE =================
\title{Uncertainty-Aware Predictive Traffic Engineering in
Software-Defined Networks: A Conformal-Gated Controller
Evaluated over 64 Controlled Runs}

% TODO: confirm co-authors, order, and affiliations before submission.
\author{\IEEEauthorblockN{Mohit Pingale}
\IEEEauthorblockA{Department of Computer Engineering\\
Vishwakarma Institute of Technology, Pune, India}}

\begin{document}
\maketitle

% ================= ABSTRACT =================
\begin{abstract}
Predictive controllers for software-defined networks (SDN) are usually
evaluated on point-forecast accuracy and a handful of runs. We argue
that neither is sufficient: a forecast that drives rerouting should
carry a measured uncertainty statement, and its value should be judged
at the level of network outcomes under repeated, configuration-matched
trials. We present a predictive traffic-engineering controller for a
Ryu/OpenFlow network that couples a gradient-boosting utilization
forecaster with split-conformal prediction intervals, an explicit
confidence gate, per-decision explanations, and drift detection, and
we evaluate four controllers (static, reactive, predictive-persistence,
predictive-ML) over 64 controlled Mininet runs (16 repetitions per
system, two campaigns, one frozen source tree, automated configuration
validation). The predictive-ML controller eliminated congestion
entirely (0.000\,s mean congestion duration and a congestion-avoidance
rate of 1.000, versus 3.000\,s and 0.250 for both reactive and
persistence control, and 38.500\,s for static routing). On 2{,}736
forecast samples paired with the identical recorded actuals, the
learned forecaster is certified better than persistence on squared
loss (Diebold--Mariano statistic 2.945, $p = 0.0032$; per-run sign
test 16/16, $p \approx 3.1 \times 10^{-5}$), while persistence remains
better on mean absolute error (10.90 vs.\ 14.78 points) and the
predictive controller pays a small, now statistically resolved
average-latency cost of $\approx$0.58\,ms. We report the negative
results with equal prominence: the conformal interval's headline
coverage is an in-sample figure whose split evaluation is pending,
the confidence gate never passed in live operation (reroutes were
executed by a measured fallback path), and a graph neural network
forecaster does not beat persistence on real telemetry. The
contribution is an uncertainty-aware control architecture and an
evaluation discipline---matched configurations, same-target forecast
tests, and intervention-aware interpretation---that makes both the
positive and the negative claims auditable.
\end{abstract}

\begin{IEEEkeywords}
software-defined networking, traffic engineering, conformal
prediction, gradient boosting, Diebold--Mariano test, OpenFlow,
Mininet
\end{IEEEkeywords}

% ================= INTRODUCTION =================
\section{Introduction}
Traffic engineering (TE) in software-defined networks increasingly
relies on short-horizon forecasts: if a controller can see congestion
five seconds before it forms, it can move traffic before queues build
rather than after loss appears. The published pattern, however, is
fragile in two ways. First, forecast quality is typically reported as
a single accuracy number on a single run, with no significance test
and no separation between the forecaster's skill and the controller's
effect on the very traffic being forecast. Second, point forecasts
give the control loop no way to distinguish a confident prediction
from a guess, so systems reroute on noise or fail to reroute on
signal.

This paper takes the position that predictive TE should be built and
judged as an uncertainty-aware system:

\begin{itemize}
\item \textbf{Intervals, not points.} The forecaster ships with
split-conformal intervals (with conformalized quantile regression and
adaptive conformal inference variants), so every routing decision can
consult a coverage-calibrated range rather than a bare number.
\item \textbf{An explicit gate, and an honest fallback.} A confidence
gate decides whether a forecast is strong enough to act on; when it
is not, a measured fallback path---not silence---handles the decision,
and every decision is logged with its reason and its top contributing
features.
\item \textbf{Evaluation that can falsify the claims.} Four
controllers are compared over 64 configuration-matched runs; forecast
comparisons pair both forecasters against identical recorded actuals
and are reported with Diebold--Mariano (DM) tests; interval coverage
is labelled by whether it was measured in-sample or held out.
\end{itemize}

The results support a narrower and stronger claim than ``ML beats
baselines.'' Predictive-ML control eliminated congestion in all 16 of
its runs, and its forecaster is significantly better than persistence
at avoiding large errors---the errors that precede congestion---while
being worse on average absolute error and slightly worse on mean
latency. Section~\ref{sec:results} states each claim with its test;
Section~\ref{sec:threats} states what the evaluation cannot support.

% ================= RELATED WORK =================
\section{Related Work}
Centralized TE with global visibility originates with systems such as
B4~\cite{jain2013b4} and SWAN~\cite{hong2013swan} over the OpenFlow
control abstraction~\cite{mckeown2008openflow}; our controller is a
single-switch-pair prototype in that lineage, not a WAN system.
Learning-based TE typically optimizes routing from predicted traffic
matrices; we instead forecast per-link utilization at a 5-second
horizon and focus on the evaluation of that forecast. Conformal
prediction~\cite{lei2018distributionfree} provides
distribution-free prediction intervals under exchangeability;
conformalized quantile regression~\cite{romano2019conformalized} and
adaptive conformal inference~\cite{gibbs2021adaptive} address
heteroscedasticity and distribution shift, both of which appear in
our data. Forecast comparison follows Diebold and
Mariano~\cite{diebold1995comparing} with Newey--West HAC variance and
the Harvey--Leybourne--Newbold small-sample correction. Graph neural
forecasters~\cite{kipf2017gcn} are a natural extension to joint
multi-link prediction; we evaluate one and report its negative
result. Experiments run in Mininet~\cite{lantz2010mininet} under the
Ryu controller framework.

% ================= SYSTEM DESIGN =================
\section{System Design}

% FIGURE 1 SOURCE: architecture of path_switch_controller.py with
% conformal_runtime.py, confidence_policy.py, explainability.py;
% redraw from frontend dashboard topology view (diamond4).
\begin{figure}[t]
\centering
\fbox{\parbox{0.92\columnwidth}{\centering\vspace{1.2em}
[Figure 1 placeholder --- system architecture: telemetry loop,
forecaster, conformal interval service, confidence gate and measured
fallback, reroute actuator, explanation and drift logs]\vspace{1.2em}}}
\caption{Controller architecture. Every reroute decision is logged
with its reason, the interval that accompanied the forecast, and the
features that contributed most.}
\label{fig:arch}
\end{figure}

\subsection{Network and workload}
The evaluation topology (Fig.~\ref{fig:arch}) is a four-switch diamond:
hosts $h_1, h_2$ attach to $s_1, s_4$; path~A ($s_1$--$s_2$--$s_4$)
has 10\,Mbps links with 2\,ms delay per link, path~B
($s_1$--$s_3$--$s_4$) 20\,Mbps with 4\,ms. A single UDP iperf3 flow
follows a fixed five-phase workload (warmup, ramp, near-threshold,
congestion, recovery; 90\,s of traffic after a 5\,s settle). The
congestion threshold is 85\% utilization and the recovery threshold
70\%, identical for every controller and run.

\subsection{Controllers under comparison}
Four systems share the same telemetry, topology, and workload and
differ only in control law: \emph{static} (never reroutes),
\emph{reactive} (reroutes on measured threshold crossings),
\emph{predictive-persistence} (predictive control whose forecaster is
$\hat{u}_{t+5} = u_t$), and \emph{predictive-ML} (predictive control
with the learned forecaster). Predictive control requires three
consecutive signals, a minimum projected improvement of 5 points, and
observes a 15\,s cooldown; transitions are make-before-break and are
timed and logged.

\subsection{Forecasters}
The learned forecaster is a scikit-learn gradient boosting regressor
(100 depth-2 estimators, learning rate 0.03, high-load samples
weighted $\times$20) over lagged utilization, rolling means, and
throughput features, trained exclusively on a separate calibration
campaign and frozen before evaluation; a portable tree representation
with checksum-verified host/guest parity is what the controller
executes. On the calibration campaign's held-out split it records
RMSE 16.11 and $R^2 = 0.700$, against linear regression (19.57,
0.558) and persistence (RMSE 22.32, MAE 9.92, $R^2 = 0.425$)
(source: \texttt{cnt-ml/results/metrics\_comparison.csv}).

\subsection{Uncertainty: conformal intervals and the gate}
Split-conformal calibration on held-out calibration-campaign rows
produces a residual radius; conformalized quantile regression wraps
the quantile forecaster (P50/P90/P99), and adaptive conformal
inference (ACI) adjusts the miscoverage level online. At decision
time the confidence gate admits a forecast-driven reroute only if the
interval's lower bound clears the congestion threshold. When the gate
does not admit the forecast---or no forecast exists---the controller
does not stall: a \emph{measured fallback} reroutes on directly
measured recovery/congestion evidence, and the decision log records
which path produced each reroute. Section~\ref{sec:gate} reports how
often each path fired in live operation; the gate's live record is
part of our results, not an assumption.

\subsection{Explanations and drift detection}
Every evaluated decision writes an explanation row (active and
candidate path, reason, gate outcome, projected improvement, top
features by split-count $\times$ deviation importance). Two drift
detectors---Page--Hinkley and an ADWIN-style window comparison---run
on the forecast-error stream and log events; an optional online
recalibration strategy adapts the conformal radius when drift is
declared. These components are evaluated as instruments: we report
what they recorded, including when they recorded nothing actionable.

\subsection{A graph forecaster, evaluated and rejected}
A two-layer graph convolutional network over the link line-graph was
trained to forecast all links jointly and compared against
persistence and a per-link ridge on identical features. Its result is
reported in Section~\ref{sec:gnn}; it is not part of the deployed
controller.

% ================= PROTOCOL =================
\section{Evaluation Protocol}
\label{sec:protocol}
\textbf{Matched configurations, frozen source.} All 64 runs (two
campaigns of 32; 16 repetitions per system in the pooled analysis) ran
one frozen source tree whose files are SHA-256 recorded in every
run's metadata. The pooled analysis re-validates, for every run, that
capacities, topology, workload, thresholds, environment versions, and
source hashes are identical, and refuses to aggregate otherwise.
Runs failing the automated quality gate (controller errors, missing
receiver data, prediction coverage below 95\% of eligible rows) are
excluded and preserved; campaign-01 was excluded in full after a
model-deserialization failure and plays no role in any number below.

\textbf{Same-target forecast tests.} A forecast comparison is only
meaningful against a fixed future. We therefore pair each forecast
issued inside the predictive-ML captures with the telemetry value
recorded for the same link exactly 5\,s later, and score
\emph{both} forecasters against that identical actual. A
cross-controller pairing---scoring one controller's forecasts against
another controller's realized traffic---is invalid, because each
controller induces a different future; an earlier pooled test of that
form was retracted for exactly this reason and is not used. DM tests
use squared loss, Newey--West HAC variance (lag 5), and the HLN
correction; a per-run sign test accompanies them.

\textbf{Interval honesty.} Coverage figures are labelled in-sample
(measured on calibration rows) or held-out (chronological split or an
untouched capture). No headline claim in this paper rests on an
in-sample coverage number; see Section~\ref{sec:conformal}.

\textbf{Intervention censoring.} When a forecast triggers a reroute,
the congestion it predicted may never occur. Counting such forecasts
as false alarms would penalize the controller for succeeding; we
therefore judge forecasts by error against realized values
(Section~\ref{sec:forecast}) and judge the \emph{system} by realized
congestion (Section~\ref{sec:congestion}), and say so explicitly
where the two readings diverge.

% ================= RESULTS =================
\section{Results}
\label{sec:results}

\subsection{Congestion: the system-level claim}
\label{sec:congestion}
Table~\ref{tab:systems} reports the pooled n=16 analysis
(source: \texttt{combined-16/final-analysis/network\_statistics.csv}).
Predictive-ML is the only controller with zero congestion: mean
congestion duration 0.000\,s (95\% CI [0.000, 0.000]) and avoidance
rate 1.000 in all 16 runs, against 3.000\,s for both reactive and
predictive-persistence control, 38.500\,s for static routing, and
avoidance 0.250 [0.112, 0.388] for both non-ML rerouting controllers.
Static routing additionally loses 11.44\% of packets; all rerouting
controllers are at or below 0.012\%.

\begin{table*}[t]
\centering
\caption{Pooled system results, 16 repetitions per system (64 runs).
Mean [95\% CI]. Source:
\texttt{final-evaluation-artifacts/combined-16/final-analysis/network\_statistics.csv}.}
\label{tab:systems}
\begin{tabular}{lrrrr}
\toprule
Metric & Static & Reactive & Pred.-Persistence & Pred.-ML \\
\midrule
Congestion duration (s)      & 38.500 [38.225, 38.775] & 3.000 [3.000, 3.000] & 3.000 [3.000, 3.000] & \textbf{0.000 [0.000, 0.000]} \\
Congestion events            & 1.500 [1.225, 1.775] & 1.000 & 1.000 & \textbf{0.000} \\
Avoidance rate               & --- & 0.250 [0.112, 0.388] & 0.250 [0.112, 0.388] & \textbf{1.000 [1.000, 1.000]} \\
Avg latency (ms)             & 21.031 [19.489, 22.572] & \textbf{11.846 [11.767, 11.925]} & 11.973 [11.766, 12.180] & 12.427 [12.215, 12.638] \\
P95 latency (ms)             & 125.969 [125.626, 126.311] & 16.676 [16.601, 16.752] & 16.734 [16.633, 16.835] & 16.773 [16.645, 16.900] \\
Avg jitter (ms)              & 0.118 [0.111, 0.126] & 0.061 [0.039, 0.083] & 0.059 [0.049, 0.069] & 0.179 [$-$0.046, 0.405] \\
Throughput (Mbps)            & 6.970 & 7.883 & 7.882 & 7.881 \\
Packet loss (\%)             & 11.442 & 0.000 & 0.012 & 0.012 \\
Mean inference (ms)          & --- & --- & 0.011 & 1.328 [1.208, 1.448] \\
Mean controller CPU (\%)     & 0.611 & 0.659 & 0.733 & 0.903 \\
\bottomrule
\end{tabular}
\end{table*}

% FIGURE 2 SOURCE: final-evaluation-artifacts/combined-16/final-analysis/network_statistics.csv
% (Congestion_Duration_seconds by system, with campaign-02-only overlay).
\begin{figure}[t]
\centering
\fbox{\parbox{0.92\columnwidth}{\centering\vspace{1.2em}
[Figure 2 placeholder --- congestion duration by controller,
n=16, with n=8 campaign-02 overlay]\vspace{1.2em}}}
\caption{Congestion duration per controller. Predictive-ML is zero
in every repetition of both campaigns.}
\label{fig:congestion}
\end{figure}

\subsection{The latency cost, stated plainly}
At n=8 the latency comparison was unresolved (overlapping intervals).
At n=16 it resolves, against the predictive controller: mean latency
12.427\,ms [12.215, 12.638] versus reactive 11.846\,ms
[11.767, 11.925]---a real cost of $\approx$0.58\,ms ($\approx$5\%).
P95 latency intervals still overlap (16.773 [16.645, 16.900] vs.\
16.676 [16.601, 16.752]), and jitter intervals overlap widely; the
predictive controller's jitter mean (0.179) is driven by high
variance across runs rather than a consistent per-run penalty. We
report the mean-latency cost as the measured price of predictive
rerouting in this workload; throughput is indistinguishable across
rerouting controllers (7.881--7.883\,Mbps).

\subsection{Forecast quality: better where it matters, worse on average}
\label{sec:forecast}
Table~\ref{tab:forecast} scores both forecasters on the 2{,}736
same-target pairs of Section~\ref{sec:protocol}. The learned
forecaster is significantly better on squared loss (DM $= 2.945$,
$p = 0.0032$; mean squared-loss differential 39.01 points in its
favour) and wins the per-run squared-loss comparison in all 16 runs
(sign test $p \approx 3.1 \times 10^{-5}$). It is simultaneously
\emph{worse} on MAE in every run (14.78 vs.\ 10.90 pooled). Both
statements are true, and the mechanism is visible in the error
structure: the learned forecaster carries a $+9.6$-point mean bias
whenever current utilization is at least 50\%---it anticipates ramps
that plateau---so it bleeds absolute error in calm periods (calm-half
MAE 13.1 vs.\ 7.4) while avoiding the large misses at transitions
that dominate squared loss and precede congestion.

\begin{table}[t]
\centering
\caption{Same-target forecast comparison, 2{,}736 pairs pooled over
16 predictive-ML runs. Source:
\texttt{combined-16/final-analysis/combined\_summary.json}
(generator: \texttt{experiments/merge\_campaigns.py}).}
\label{tab:forecast}
\begin{tabular}{lrr}
\toprule
 & Persistence & Gradient Boosting \\
\midrule
MAE (points)  & \textbf{10.90} & 14.78 \\
RMSE (points) & 22.99 & \textbf{22.12} \\
\midrule
\multicolumn{3}{l}{DM (squared loss, HAC + HLN): 2.945, $p = 0.0032$} \\
\multicolumn{3}{l}{Per-run sign test: GB wins 16/16, $p \approx 3.1 \times 10^{-5}$} \\
\bottomrule
\end{tabular}
\end{table}

We attempted to close the MAE gap without touching the frozen model,
selecting every candidate leave-one-run-out (fit on 7 runs, scored on
the held-out run, campaign-02's 8 runs): a volatility gate, a
predicted-move gate, and convex blends all collapse to pure
persistence in every fold; linear recalibration attains the best RMSE
of any variant (21.25) but worsens MAE (15.67). An oracle gate
(MAE 8.43) shows headroom exists, but no signal available at
prediction time identifies when the learned forecaster is right. We
therefore make no hybrid claim; the MAE gap is a training-distribution
artifact whose honest remedy---retraining on calibration data only,
never on the evaluation campaigns---is future work.

% FIGURE 3 SOURCE: paired forecast errors from
% final-evaluation-artifacts/campaign-02 + campaign-03 predictive-ml
% predictions.csv/telemetry.csv (audit scripts: WEAK_METRIC_AUDIT.md).
\begin{figure}[t]
\centering
\fbox{\parbox{0.92\columnwidth}{\centering\vspace{1.2em}
[Figure 3 placeholder --- forecast error by regime: calm vs.\
volatile halves, MAE and RMSE for both forecasters]\vspace{1.2em}}}
\caption{The learned forecaster's MAE deficit is concentrated in
calm periods; its squared-loss advantage is concentrated at
transitions.}
\label{fig:regime}
\end{figure}

\subsection{Conformal intervals: an in-sample headline, a pending split}
\label{sec:conformal}
The shipped calibration artifact
(\texttt{cnt-ml/models/conformal\_calibration.json}; calibrated on
172 rows of one campaign-02 capture) carries radius 38.51 points and
an empirical coverage of 0.907 at nominal 0.90. \textbf{That coverage
figure is in-sample}---it is measured on the calibration rows
themselves---and we do not claim it as held-out evidence. A
chronological split evaluation of the interval (calibrate on an early
segment, measure on a later, untouched segment and on an untouched
capture) is the required correction and is in progress; its result
belongs here: \textcolor{red}{\textbf{[SPLIT COVERAGE: TBD]}}.
Until it lands, the interval's status in this paper is a mechanism
with a measured calibration, not a certified guarantee---and the
gate behaviour reported next is exactly what that status predicts.

\subsection{The confidence gate never passed live; the fallback operated}
\label{sec:gate}
Across all recorded live captures, the confidence gate admitted zero
forecast-driven reroutes (690 recorded live intervals, 0 passed). In
the fully instrumented predictive-safe live run, all 140 intervals
were held by the gate; the run's two reroutes were executed by the
measured fallback, and its 80 logged decisions decompose as 40
\texttt{confidence\_gate\_closed}, 30 \texttt{measured\_recovery},
and 10 \texttt{missing\_prediction}, with utilization the top
contributing feature. We document this rather than tune it away:
with intervals as wide as the current calibration produces, a gate
that requires the lower bound to clear an 85\% threshold is close to
arithmetically dormant, and shrinking intervals or lowering the
threshold to manufacture passes would destroy the meaning of both.
The system remained safe and effective because the fallback is a
first-class path, and the explanation log makes the substitution
visible decision by decision. In the same live run, both drift
detectors (Page--Hinkley and ADWIN-style) fired on the error stream,
and 66 anomalies (62 elephant-flow, 4 surge) were recorded; live
transitions completed in 0.534\,s on average with no failures.

\subsection{The GNN forecaster does not beat persistence}
\label{sec:gnn}
On real campaign telemetry the graph forecaster records MAE 18.39
against persistence 14.87 (RMSE 27.37 vs.\ 26.11; DM $p = 0.491$, not
significant), after fixes for exact-timestamp windowing and residual
self-projections. The structural cause is visible in the data: the
legacy captures record only the two $s_1$ uplinks, so the line-graph
has two nodes and the joint model has almost no spatial structure to
exploit. The controller now records every core link
(\texttt{telemetry\_all\_links.csv}); the pooled re-evaluation on
that data is the correct test and has not yet been run. The negative
result stands as reported.

% ================= THREATS =================
\section{Threats to Validity}
\label{sec:threats}
\textbf{Single topology and workload family.} All controlled results
come from one diamond topology and one five-phase workload; the
access6 topology is exercised only in non-controlled scenarios.
Generalization is a claim we do not make.
\textbf{Emulation.} Mininet timings are not hardware timings;
absolute latencies would differ on physical switches, though the
matched design protects the comparisons.
\textbf{Exchangeability.} Conformal coverage is marginal and assumes
exchangeability within the calibration window; the workload's phase
structure violates it in principle, which is why the split evaluation
of Section~\ref{sec:conformal}---not the in-sample figure---is the
number that matters, and why ACI is included as mitigation rather
than proof.
\textbf{Intervention censoring.} Forecast errors are scored against
futures the controller itself altered; Section~\ref{sec:protocol}
describes the pairing discipline this forced, and the retracted
cross-controller test is the cautionary example.
\textbf{One-second sampling} may miss sub-second congestion, and
controller event timestamps do not certify hardware flow
acknowledgement.

% ================= CONCLUSION =================
\section{Conclusion}
A predictive SDN controller earns its complexity only if its
advantage survives repetition, matched configuration, and tests that
could have gone the other way. Over 64 controlled runs, ours does on
the outcome it exists for---congestion is eliminated in every
repetition---and its forecaster is certified better than persistence
precisely on the large errors that precede congestion, while costing
0.58\,ms of mean latency and losing on average absolute error. Its
uncertainty machinery is reported at the same standard: an in-sample
coverage headline awaiting its split evaluation, and a confidence
gate that never passed live, with the measured fallback---fully
logged---doing the work instead. We regard that combination, not any
single metric, as the publishable unit: an uncertainty-aware
controller whose strengths and failure modes are both measured.

\section*{Artifact Availability}
All code, run captures, campaign manifests, and analysis scripts are
public at
\url{https://github.com/mohitpingale24-alt/Predictive-SDN-Traffic-Engineering}.
The pooled statistics in Table~\ref{tab:systems} regenerate from the
preserved captures with \texttt{experiments/merge\_campaigns.py};
every run's metadata records the SHA-256 of the source files that
produced it.

% ================= REFERENCES =================
\begin{thebibliography}{00}
\bibitem{mckeown2008openflow} N.~McKeown \emph{et al.}, ``OpenFlow:
Enabling innovation in campus networks,'' \emph{ACM SIGCOMM CCR},
vol.~38, no.~2, pp.~69--74, 2008.
\bibitem{jain2013b4} S.~Jain \emph{et al.}, ``B4: Experience with a
globally-deployed software defined WAN,'' in \emph{Proc. ACM
SIGCOMM}, 2013, pp.~3--14.
\bibitem{hong2013swan} C.-Y.~Hong \emph{et al.}, ``Achieving high
utilization with software-driven WAN,'' in \emph{Proc. ACM SIGCOMM},
2013, pp.~15--26.
\bibitem{lantz2010mininet} B.~Lantz, B.~Heller, and N.~McKeown, ``A
network in a laptop: Rapid prototyping for software-defined
networks,'' in \emph{Proc. ACM HotNets}, 2010.
\bibitem{lei2018distributionfree} J.~Lei, M.~G'Sell, A.~Rinaldo,
R.~J.~Tibshirani, and L.~Wasserman, ``Distribution-free predictive
inference for regression,'' \emph{J. Amer. Statist. Assoc.},
vol.~113, no.~523, pp.~1094--1111, 2018.
\bibitem{romano2019conformalized} Y.~Romano, E.~Patterson, and
E.~J.~Cand\`es, ``Conformalized quantile regression,'' in
\emph{Proc. NeurIPS}, 2019.
\bibitem{gibbs2021adaptive} I.~Gibbs and E.~J.~Cand\`es, ``Adaptive
conformal inference under distribution shift,'' in \emph{Proc.
NeurIPS}, 2021.
\bibitem{diebold1995comparing} F.~X.~Diebold and R.~S.~Mariano,
``Comparing predictive accuracy,'' \emph{J. Bus. Econ. Statist.},
vol.~13, no.~3, pp.~253--258, 1995.
\bibitem{kipf2017gcn} T.~N.~Kipf and M.~Welling, ``Semi-supervised
classification with graph convolutional networks,'' in \emph{Proc.
ICLR}, 2017.
\end{thebibliography}

\balance
\end{document}
